% from ega0-21-fr.tex
\oldpage[0\textsubscript{IV-1}]{169}
cohomologie, et notons que, en vertu de l'exactitude de (21.6.1.1), la
cohomologie de $T_1$ et celle de $T_2$ sont nulles sauf en un seul et même
degré, pour lequel les modules de cohomologie sont tous deux égaux à
$\Upsilon_{k''/k'/k}\otimes_{k''}L$~; comme $T_1$ et $T_2$ ont ainsi même
cohomologie, celle de $T_3$ est nécessairement nulle, ce qui prouve (i).
  \item[(ii)] Considérons de même le diagramme commutatif
\[
\xymatrix{
0\ar[r]&\Delta(L/K,k/k_0)\otimes_LM\ar[r]\ar[d]&
\Delta(M/K,k/k_0)\ar[r]\ar[d]&\Delta(M/L,k/k_0)\ar[r]&0\\
0\ar[r]&\Upsilon_{L/K/k_0}\otimes_LM\ar[r]\ar[d]&
\Upsilon_{M/K/k_0}\ar[r]\ar[d]&\Upsilon_{M/L/k_0}\ar[r]&0\\
0\ar[r]&\Upsilon_{L/K/k}\otimes_LM\ar[r]&
\Upsilon_{M/K/k}\ar[r]&\Upsilon_{M/L/k}\ar[r]&0
}
\]
où de nouveau on considère les trois lignes comme des \emph{complexes}
$T'_1,T'_2,T'_3$~; la suite exacte (21.6.1.1) et la définition de $\Delta$
(21.6.1.2) donnent ici une suite exacte de complexes
$0\to T'_1\to T'_2\to T'_3\to0$ à laquelle on applique encore la suite exacte
de cohomologie~; cette fois, en vertu de l'exactitude de (21.6.1.1), la
cohomologie de $T'_2$ et celle de $T'_3$ sont nulles sauf en un seul et même
degré, pour lequel les modules de cohomologie sont tous deux égaux à
$\Upsilon_{M/L/k_0}$~; on conclut donc ici que la cohomologie de $T'_1$ est
nulle, ce qui établit (ii).
\end{enumerate}
\end{proof}

\begin{corollary}[21.6.5]
\label{0.21.6.5-fr}
\begin{enumerate}
  \item[(i)] Étant donnés cinq corps $k\subset k'\subset k''\subset K\subset L$,
  pour que $k''$ soit $k$-admissible pour l'extension $L$ de $K$, il faut et
  il suffit que $k'$ soit $k$-admissible et que $k''$ soit $k'$-admissible
  pour l'extension $L$ de $K$.
  \item[(ii)] Étant donnés cinq corps
  $k_0\subset k\subset K\subset L\subset M$, pour que $k$ soit
  $k_0$-admissible pour l'extension $M$ de $K$, il faut et il suffit qu'il le
  soit pour l'extension $L$ de $K$ et pour l'extension $M$ de $L$.
\end{enumerate}

Cela résulte aussitôt des relations (21.6.4.2) et (21.6.4.4), les valeurs de
$d$ étant $\geqslant0$.
\end{corollary}

\begin{corollary}[21.6.6]
\label{0.21.6.6-fr}
Soient $k_0\subset k\subset K\subset L$ quatre corps, et supposons que $k$
soit $k_0$-admissible pour l'extension $L$ de $K$. Alors, si $k'_0$, $k'$,
$K'$, $L'$ sont quatre corps tels que
$k_0\subset k'_0\subset k'\subset k\subset K\subset K'\subset L'\subset L$,
$k'$ est $k'_0$-admissible pour l'extension $L'$ de $K'$.
\end{corollary}

\subsection{L'égalité de Cartier}
\label{subsection:0.21.7-fr}

Le résultat suivant traduit en termes de différentielles un théorème de
MacLane sur les dérivations~:

